\documentclass{article}
\usepackage{amsmath}
\usepackage{amsthm}

\newtheorem{theorem}{Problem}

\begin{document}

\begin{theorem}
Compute $58_9 - 18_9.$ Express your answer in base $9.$
\end{theorem}

\begin{proof}
Step 0: Step 1. I need to subtract two numbers in base 9: 58₉ minus 18₉. I'm supposed to express the answer in base 9 as well.
Step 1: Step 2. First, I should recall how subtraction works in different bases. It's similar to subtraction in base 10, but I have to borrow according to the base I'm working in.
Step 2: Step 3. Let's write the numbers one above the other, aligning the digits by place value:

                 5 8₉

               - 1 8₉

               --------
Step 3: Step 4. Starting from the rightmost digit, I need to subtract 8₉ from 8₉.
Step 4: Step 5. 8 minus 8 is 0, so the rightmost digit of the answer is 0.
Step 5: Step 6. Now, moving to the next digit to the left, I need to subtract 1₉ from 5₉.
Step 6: Step 7. 5 minus 1 is 4, so the next digit is 4.
Step 7: Step 8. So, putting it all together, 58₉ minus 18₉ should be 40₉.
Step 8: Step 9. Wait a minute, that seems straightforward, but I should double-check if there's any borrowing involved.
Step 9: Step 10. In this case, since 8 - 8 = 0 and 5 - 1 = 4, there's no need for borrowing.
Step 10: Step 11. To confirm, I can convert both numbers to base 10, perform the subtraction in base 10, and then convert the result back to base 9.
Step 11: Step 12. Converting 58₉ to base 10: 5*9^1 + 8*9^0 = 5*9 + 8*1 = 45 + 8 = 53.
Step 12: Step 13. Converting 18₉ to base 10: 1*9^1 + 8*9^0 = 1*9 + 8*1 = 9 + 8 = 17.
Step 13: Step 14. Subtracting in base 10: 53 - 17 = 36.
Step 14: Step 15. Now, converting 36 back to base 9.
Step 15: Step 16. To find the digits in base 9, I divide 36 by 9, which gives 4 with a remainder of 0.
Step 16: Step 17. So, 36 in base 10 is 40 in base 9, which matches my earlier result.
Step 17: Step 18. Therefore, the answer is indeed 40₉.

        # Final Answer

        \boxed{40_9}
\end{proof}

\end{document}
